% from main.tex
%%%%%%%%%%  scan idx 479  |  folio 466  %%%%%%%%%%
\section*{Exposé VIII. Le théorème de Riemann--Roch}
\begin{center}
par Pierre Berthelot
\end{center}

\subsection*{1. Morphismes d'intersection complète}

\textbf{Définition 1.1.} Soit $f:X\to Y$ un morphisme de schémas. On dit que $f$ est un \emph{morphisme d'intersection complète} si pour tout point $x$ de $X$, il existe un voisinage ouvert $U$ de $x$, un $Y$-schéma lisse $V$, tels que la restriction de $f$ à $U$ se factorise en
\[
\begin{tikzcd}[column sep=tiny]
U \arrow[r,"i"] \arrow[dr] & V \arrow[d]\\
& Y,
\end{tikzcd}
\]
où $i$ est une immersion régulière (VII 1.4). Quitte à remplacer $V$ par un ouvert, on peut supposer que $i$ est une immersion fermée régulière.

On peut dans cette définition remplacer $V$ par un fibré vectoriel $Y[T_1,\ldots,T_n]$. En effet, la notion étant locale, on peut supposer que $V$ est affine, et que son image dans $Y$ est contenue dans un ouvert affine $W$ de $Y$. Il existe alors un entier $n$ et une $W$-immersion fermée $V\longrightarrow W[T_1,\ldots,T_n]$, d'où une $Y$-immersion $V\to Y[T_1,\ldots,T_n]$, qui est régulière d'après VII 1.10. Comme la composée de deux immersions régulières est régulière, on trouve donc une $Y$-immersion régulière $U\longrightarrow Y[T_1,\ldots,T_n]$ factorisant $f$.

\textbf{Proposition 1.2.} Soit $f:X\to Y$ un morphisme d'intersection complète. Alors toute $Y$-immersion $i:X\to Z$ de $X$ dans un $Y$-schéma lisse est régulière.

Nous montrerons d'abord le corollaire :

\textbf{Corollaire 1.3.} Soit le diagramme commutatif :
\[
\begin{tikzcd}[column sep=large,row sep=large]
& V' \arrow[d,"p"]\\
X \arrow[ur,"j"] \arrow[r,"i"'] & V
\end{tikzcd}
\]
